\documentclass[10pt,a4paper]{amsart}
\usepackage[margin=19mm]{geometry}
\usepackage{amsmath,amssymb,mathtools}
\usepackage[scr=rsfs,cal=euler]{mathalfa}
\usepackage{xfrac}
\usepackage[unicode,hidelinks,bookmarks=false]{hyperref}
\newcommand{\ie}{i.e.\ }
\newcommand{\cf}{cf.\ }
\newcommand{\resp}{respectively\ }
\newcommand{\Subsection}{\subsection}
\providecommand{\nobreakdash}{\nobreak\hbox{-}}
\setlength{\emergencystretch}{2em}
\multlinegap0pt
\usepackage{bm}
\usepackage{bbm}
\usepackage{dsfont}
\usepackage{xspace}
\usepackage{cancel}
\usepackage{mathtools}
\usepackage{amsthm}
\makeatletter
\@ifundefined{theorem}{\newtheorem{theorem}{Theorem}}{}
\makeatother
\title[Maximal Achievable Service Rates of Codes and Connections to]{Maximal Achievable Service Rates of Codes and Connections to Combinatorial Designs}
\author{Hoang Ly, Emina Soljanin}
\date{Source: arXiv:2506.16983v3 (CC BY 4.0)}
\begin{document}
\maketitle

\noindent\emph{Source and licence.} Maximal Achievable Service Rates of Codes and Connections to Combinatorial Designs. Version arXiv:2506.16983v3 (CC BY 4.0), distributed under the Creative Commons Attribution 4.0 International licence, as recorded in the arXiv licence metadata for this exact version and shown on the abstract page https://arxiv.org/abs/2506.16983v3. Full rights notice: licenses/arxiv-2506.16983-CC-BY-4.0.txt.\par\smallskip

\noindent\textit{Editorial scope.} Counted unit: the Theorem whose source label is \texttt{thm:general\_bound\_demand} in the article (source0.tex lines 373--379, source label \texttt{thm:general\_bound\_demand}), delivered together with the article's own complete original proof of it (source0.tex lines 381--440, one \texttt{proof} environment of 60 lines). No mathematics is written, repaired or re-derived: statement and proof are the article's own text line for line. Required context retained uncounted: none. Every definition, display and label of those lines is retained unchanged. Omitted from this package and not redistributed: 1--372, 380--380, 441--925. Nothing inside a retained excerpt is omitted or altered: every retained body is delivered exactly as the article's lines are written, so no omission receipt of any kind accompanies this package. Citations: the retained excerpts carry no citation command at all, so no bibliography entry of the article is transcribed and no reference is redistributed. Reference adaptation: none was needed and none was made; every reference inside the delivered unit resolves inside it, so no unresolved reference or citation is produced. Printed slips kept literal: no printed word, symbol, hypothesis or number of the counted statement or of its proof was altered. Layout and class: the article's own class is \texttt{IEEEtran} with packages including balance, inputenc, fontenc, url, ifthen, cite, enumitem, amsmath, graphicx, mathtools, amssymb, amsthm, verbatim, hyperref; the wrapper is the lane's pinned \texttt{amsart} editorial wrapper at 10pt on A4, which restates the theorem-like environments the retained text needs and adds the article's own macro definitions that the retained text needs (none), with extra wrapper packages (bm, bbm, dsfont, xspace, cancel, mathtools). The delivered numbering is the unit's own. Assets: no retained span contains a figure environment, a graphics-inclusion command, a PDF-page inclusion, a table or a TikZ picture; no image file is copied into this package, the package declares no asset dependency and no image file is redistributed with it. Licence: version arXiv:2506.16983v3 of this article is distributed under the Creative Commons Attribution 4.0 International licence, as recorded in the arXiv licence metadata for this exact version and shown on the abstract page https://arxiv.org/abs/2506.16983v3. A full rights notice is delivered at licenses/arxiv-2506.16983-CC-BY-4.0.txt. Characters: every retained excerpt is pure ASCII, so the delivered engine, XeLaTeX with its default Latin Modern text font, reports no missing character.
\begin{theorem}\label{thm:general_bound_demand}
Consider a storage system encoded by a generator matrix \( \boldsymbol{G} \). Let \( \mathcal{O} \) be a smallest recovery set for object \( o_\ell \), with \( \ell \in [k] \), and let \( a = |\mathcal{O}| \) denote its size. Let \( J_{\mathcal{O}} \) be the maximum number of parity-check sums orthogonal on \( \mathcal{O} \). Then the maximal achevable rate \( \lambda_\ell^{\max} \) for object \( o_\ell \) satisfies
\[
1 + J_{\mathcal{O}} \,\le\, \lambda_{\ell}^{\max} \,\le\, 1 + \frac{n - a}{\max\{d^\perp - a,\ a\}} \,\le\, 1 + \frac{n - a}{d^\perp - a},
\]
where \( d^{\perp} \) is the minimum distance of the dual code \( \mathcal{C}^{\perp} \).
\end{theorem}

\begin{proof}
The lower follows from the above argument. We now prove the upper bound.
\\[2ex]
\textbf{Step 1: Every recovery set other than \( \mathcal{O} \) has size at least \( d^\perp - a \).}

Let \( \mathcal{O} \) be a minimal recovery set for \( o_\ell \), with size \( a = |\mathcal{O}| \), and let \( \mathcal{O}' \ne \mathcal{O} \) be any other recovery set for \( o_\ell \). Since both \( \mathcal{O} \) and \( \mathcal{O}' \) recover \( o_{\ell} \), we have
\[
\sum_{j \in \mathcal{O}} \boldsymbol{c}^j \,=\, \boldsymbol{e}_{\ell} \,=\, \sum_{j \in \mathcal{O}'} \boldsymbol{c}^j.
\]
Adding these equations over \( \mathbb{F}_2 \) yields
\[
\sum_{j \in \mathcal{O}} \boldsymbol{c}^j + \sum_{j \in \mathcal{O}'} \boldsymbol{c}^j \,=\, \boldsymbol{0}_k.
\]
Let \( \mathcal{O}_1 := (\mathcal{O} \cup \mathcal{O}') \setminus (\mathcal{O} \cap \mathcal{O}') \) be the symmetric difference of \( \mathcal{O} \) and \( \mathcal{O}' \). Then
\[
\sum_{j \in \mathcal{O}_1} \boldsymbol{c}^j \,=\, \boldsymbol{0}_k,
\]
which implies that the incidence vector \( \boldsymbol{\chi}(\mathcal{O}_1) \in \mathbb{F}_2^n \) satisfies
\[
\boldsymbol{G} \cdot \boldsymbol{\chi}(\mathcal{O}_1)^\top \,=\, \boldsymbol{0}_k,
\]
so \( \boldsymbol{\chi}(\mathcal{O}_1) \in \mathcal{C}^\perp \). Since \( \mathcal{O} \ne \mathcal{O}' \), the set \( \mathcal{O}_1 \) is nonempty, and thus \( \boldsymbol{\chi}(\mathcal{O}_1) \) is a nonzero codeword in \( \mathcal{C}^\perp \) with weight at least \( d^\perp \). One has
\[
|\mathcal{O}_1| \,=\, \textsf{w}(\boldsymbol{\chi}(\mathcal{O}_1)) \,\ge\, d^\perp.
\]
By the identity
\[
|\mathcal{O}_1| \,=\, |\mathcal{O}| + |\mathcal{O}'| - 2|\mathcal{O} \cap \mathcal{O}'|,
\]
we obtain
\[
|\mathcal{O}'| \,\ge\, |\mathcal{O}_1| - |\mathcal{O}| \,\ge\, d^\perp - a.
\]
% Hence, every recovery set for \( o_{\ell} \) other than \( \mathcal{O} \) must have size at least \( d^\perp - a \). Recall that $\mathcal{O}$ is also the smallest recovery set of size $a$, we conclude that all recovery sets other than \( \mathcal{O} \) must have size at least $\max\{d^{\perp} - a,\, a\}$.
Therefore, every recovery set distinct from \( \mathcal{O} \) must have size at least \( d^\perp - a \). Since \( \mathcal{O} \) has size \( a \) and is by definition the smallest recovery set, it follows that the size of any other recovery set is at least \( \max\{d^\perp - a,\, a\} \).
\\[2ex]
\textbf{Step 2: Upper bound via capacity constraint.}
To support a demand rate \( \lambda_{\ell} = \lambda_{\ell}^{\max} \), the system must allocate server capacity across recovery sets for object \( o_{\ell} \). As shown in Step 1, the object can be recovered from
\begin{itemize}
    \item one recovery set \( \mathcal{O} \) of size \( a \), and
    \item other recovery sets, each of size at least \( \max\{d^{\perp} - a,\, a\} \ge a \).
\end{itemize}

To minimize total server usage, we consider the following assignment: allocate one unit of request to the smallest recovery set \( \mathcal{O} \), which consumes \( a \) servers; assign the remaining \( \lambda_{\ell} - 1 \) units of request to other recovery sets, each of size at least \( \max\{d^\perp - a,\, a\} \). Therefore, the total number of servers involved is at least
\[
1\cdot a + (\lambda_{\ell} - 1)\cdot\max\{d^\perp - a,\ a\}.
\]

Since this request is achievable, the total required server capacity must not exceed the total available capacity \( n \). Thus,
\[
a + (\lambda_{\ell} - 1)\cdot\max\{d^\perp - a,\ a\} \,\le\, n.
\]

Solving for \( \lambda_{\ell} \) yields
\[
\lambda_{\ell} \,\le\, 1 + \frac{n - a}{\max\{d^\perp - a,\, a\}} \,\le\, 1 + \frac{n - a}{d^\perp - a}.
\]

Since this holds for all achievable \( \lambda_{\ell} \), it applies in particular to \( \lambda_{\ell}^{\max} \). Together with the lower bound \( \lambda_{\ell}^{\max} \ge 1 + J_{\mathcal{O}} \), this completes the proof.
\end{proof}

\end{document}
